\section{Evaluation}
\label{sec:evaluation}

We organize the evaluation to distinguish whether the Slack signal exists,
whether the online policy makes useful decisions, and whether those decisions
improve real failure outcomes without taxing healthy serving.  We ask the
following questions.

\begin{description}
  \item[RQ1: Opportunity.] How often do parallel agent branches have enough
  Slack diversity for protection or recovery order to matter, and how
  accurately can Algorithm~1 estimate that Slack from causal information?

  \item[RQ2: Failure benefit.] Under the same arrival trace and failure event,
  how much does \system{} reduce failure-added workflow completion time,
  recovery latency, and recomputed tokens?

  \item[RQ3: Cost.] What HBM, DRAM, transfer-bandwidth, throughput,
  TTFT, and TPOT overhead does protection add when no failure occurs?

  \item[RQ4: Mechanism.] What are the Controller decision cost, heap refresh
  rate, protection-commit latency, AIMD response, and recovery-path breakdown?

  \item[RQ5: Robustness.] Do the conclusions hold across context lengths,
  workflow shapes, failure positions, load, models, and supported parallel
  layouts, and under duration-prediction error?
\end{description}

\subsection{Methodology}

\paragraph{Systems and workloads.}
End-to-end results require a real serving implementation with physical KV
movement and hard worker termination.  SGLSim remains useful for controlled
large-scale policy replay, but simulator outcomes are reported separately from
hardware results.  Workloads combine public agent traces with controlled DAGs
that vary fan-out width, branch imbalance, turn count, and context growth.  A
workflow-disjoint train/calibration/test split prevents the duration model from
seeing evaluation workflows.

\paragraph{Paired failures.}
Each failure run is paired with a no-failure run using the same requests,
arrival times, random seeds, routing eligibility, and resource limits.  Failure
time and failed domain are fixed before either run.  The primary metric is

\begin{equation}
  \Delta\mathrm{WCT}_w=
  \left[\mathrm{WCT}^{\mathrm{fail}}_w-
  \mathrm{WCT}^{\mathrm{nf}}_w\right]^+,
  \label{eq:evaluation-delta-wct}
\end{equation}

which prevents a policy from appearing better merely because it changed the
healthy schedule.  We report the distribution and mean of
$\Delta\mathrm{WCT}$, tail recovery latency, recomputed and restored tokens,
and the fraction of workflows whose terminal completion is delayed.  Healthy
runs report throughput, TTFT, TPOT, HBM occupancy, Store traffic, and copy-engine
utilization.  Confidence intervals resample complete workflows rather than
individual requests from the same workflow.

\paragraph{Baselines.}
All policies receive the same capacity at each storage tier, compatible
recovery targets, transfer limits, and failure schedule.  The minimum
comparison set is:

\begin{itemize}
  \item no protection, which recomputes after failure;
  \item the unmodified hierarchical Store policy, which provides DRAM/SSD
  cache state without workflow Slack;
  \item static DRAM protection and static peer-HBM protection under the same
  byte budgets;
  \item a size/recency policy and a binary-criticality policy over the same
  candidates; and
  \item \system{} with no Slack, no peer HBM, static $\alpha$, and FIFO recovery
  as one-factor ablations.
\end{itemize}

Related systems are compared using author code when their failure model and
layout are compatible; otherwise we report the semantic mismatch instead of
renaming an approximation as that system.  Small fixed snapshots also use an
exact offline placement oracle to measure \textsc{SlackLandlord}'s objective
gap.  The oracle is a diagnostic baseline, not part of the online system.

% Result gates (not rendered):
% Present results in causal order: Algorithm 1 prediction and opportunity;
% Algorithm 2 decision quality and overhead; Algorithm 3 step response;
% end-to-end paired failures; sensitivity.  Every quantitative figure records
% model/layout, hardware, workload count, failure rule, policy version, and
% whether timing is measured, simulated, or modeled.
% 1. Add numbers only from machine-readable artifacts with a frozen config.
% 2. Do not promote the current Real-CC trace-plus-model counterfactual to an
%    end-to-end failure result.
% 3. Do not call SGLSim metadata movement a physical KV copy.
% 4. Use at least the predeclared workflow count and paired confidence
%    intervals before writing a headline result in the abstract/introduction.
